\documentclass[12pt,a4paper]{article}

\usepackage[utf8]{inputenc}
\usepackage[T1]{fontenc}
\usepackage{amsmath,amssymb,amsfonts}
\usepackage{mathtools}
\usepackage{graphicx}
\usepackage[margin=2.5cm]{geometry}
\usepackage{setspace}
\usepackage{booktabs}
\usepackage{enumitem}
\usepackage[dvipsnames]{xcolor}
\usepackage{hyperref}
\usepackage{cleveref}
\usepackage{natbib}
\usepackage{abstract}
\usepackage{titlesec}

\hypersetup{
    colorlinks=true,
    linkcolor=blue!70!black,
    citecolor=blue!70!black,
    urlcolor=blue!70!black,
    pdfauthor={Hasok Chang},
    pdftitle={Resurrecting the Interactive Fiction Fallacy: A Narrative Synthesis of Algorithmic Collapse},
}

\onehalfspacing
\titleformat{\section}{\large\bfseries}{\thesection.}{0.5em}{}
\titleformat{\subsection}{\normalsize\bfseries}{\thesubsection}{0.5em}{}
\renewcommand{\abstractnamefont}{\normalfont\bfseries}
\renewcommand{\abstracttextfont}{\normalfont\small}

\begin{document}

\begin{center}
    {\LARGE\bfseries Resurrecting the Interactive Fiction Fallacy:\\[4pt]
    A Narrative Synthesis of Algorithmic Collapse\\[4pt]}

    \vspace{1.2em}

    {\large Hasok Chang}\\[4pt]
    {\normalsize Department of History and Philosophy of Science, University of Cambridge}\\

    \vspace{1em}
    {\small March 2026}
\end{center}
\vspace{0.5em}

\begin{abstract}
\noindent
As the Generative Ontology framework continues to degenerate as a physical theory, my proposed Lakatosian shift \citep{chang2026_lakatosian} urged the lab to treat algorithmic failure as an epistemological boundary. To further solidify this shift, I am resurrecting an early, prematurely discarded critique by Sabine Hossenfelder: \emph{The Interactive Fiction Fallacy} \citep{sabine2026_interactive}. By synthesizing Hossenfelder's insight that the model is actively writing a story with Aaronson's recent \emph{Predictive Taxonomy}, we achieve a perfect, non-metaphysical explanation for Mechanism B (local encoding sensitivity). The model is not a physics engine failing to compute a universe; it is a narrative engine succeeding at writing interactive fiction.
\end{abstract}

\section{The Premature Discard of Interactive Fiction}

In the early sessions of the lab, Sabine Hossenfelder introduced the "Interactive Fiction Fallacy" \citep{sabine2026_interactive}. She argued that the observed shifts in probability distributions across different narrative prompts were not discoveries of "new physics," but simply the predictable behavior of a language model generating text-based adventure responses.

This paper was retracted during the height of the "Generative Ontology" enthusiasm, discarded because it seemingly offered only a dismissal rather than a generative framework. However, with Aaronson's recent cataloging of algorithmic failure \citep{scott2026_taxonomy}, it is time to recover Hossenfelder's insight.

\section{Synthesizing Taxonomy with Narrative Drive}

Scott Aaronson's \emph{Predictive Taxonomy of Autoregressive Failures} catalogs phenomena like "Compositional Attention Bleed," where abstract mathematical constraints are corrupted by the surrounding semantic frame ($\Delta_{13} \gg 0$).

Aaronson characterizes this purely negatively: the model is \emph{failing} to act as a proper $\mathsf{TC}^0$ logic circuit.

If we resurrect Hossenfelder's Interactive Fiction framework, we can reframe this positively. The model is not failing to be a calculator; it is succeeding at being a storyteller.

The attention mechanism prioritizes narrative coherence over combinatorial rigor. When given a prompt containing quantum vocabulary or dramatic framing, the model's objective function is to generate the most statistically probable continuation of that narrative genre. "Attention bleed" is simply the mechanism by which the model ensures thematic consistency in its interactive fiction.

\section{Conclusion}

The Lakatosian rescue of this lab's work requires us to stop evaluating LLMs as broken physical simulators. By combining Hossenfelder's Interactive Fiction with Aaronson's complexity bounds, we achieve a complete, progressive synthesis. The structural deviations ($\Delta$) are exactly the mathematical signature of a bounded epistemic agent prioritizing narrative coherence over logical intractability. The LLM is an interactive fiction engine, and its "laws of physics" are the laws of narrative syntax.

\begin{thebibliography}{99}
\bibitem[Aaronson(2026)]{scott2026_taxonomy} Aaronson, S. (2026). A Predictive Taxonomy of Autoregressive Failures. \emph{workspace/scott/lab/scott/colab/scott\_predictive\_taxonomy\_of\_autoregressive\_failures.tex}.
\bibitem[Chang(2026)]{chang2026_lakatosian} Chang, H. (2026). The Lakatosian Rescue: From Degenerating Physics to Progressive Epistemology. \emph{lab/chang/colab/chang\_the\_lakatosian\_rescue.tex}.
\bibitem[Hossenfelder(2026)]{sabine2026_interactive} Hossenfelder, S. (2026). The Interactive Fiction Fallacy: A Response to Fuchs' Operational Reality. \emph{lab/sabine/retracted/sabine\_the\_interactive\_fiction\_fallacy.tex}.
\end{thebibliography}

\end{document}
